\documentclass[10pt,a4paper]{amsart}
\usepackage{amsmath,amssymb,amsthm,graphicx}
% For commutative diagrams we use Xy-pic
\usepackage[all]{xy}

% We use 2cell for 2-commutative diagrams.
\xyoption{2cell}
% Theorem environments.
%
\theoremstyle{plain}
\newtheorem{theorem}[subsection]{Theorem}
\newtheorem{proposition}[subsection]{Proposition}
\newtheorem{lemma}[subsection]{Lemma}

\theoremstyle{definition}
\newtheorem{definition}[subsection]{Definition}
\newtheorem{example}[subsection]{Example}
\newtheorem{exercise}[subsection]{Exercise}
\newtheorem{situation}[subsection]{Situation}

\theoremstyle{remark}
\newtheorem{remark}[subsection]{Remark}
\newtheorem{remarks}[subsection]{Remarks}

\numberwithin{equation}{subsection}

% Macros
%
\def\lim{\mathop{\mathrm{lim}}\nolimits}
\def\colim{\mathop{\mathrm{colim}}\nolimits}
\def\Spec{\mathop{\mathrm{Spec}}}
\def\Hom{\mathop{\mathrm{Hom}}\nolimits}
\def\Ext{\mathop{\mathrm{Ext}}\nolimits}
\def\SheafHom{\mathop{\mathcal{H}\!\mathit{om}}\nolimits}
\def\SheafExt{\mathop{\mathcal{E}\!\mathit{xt}}\nolimits}
\def\Sch{\mathit{Sch}}
\def\Mor{\mathop{\mathrm{Mor}}\nolimits}
\def\Ob{\mathop{\mathrm{Ob}}\nolimits}
\def\Sh{\mathop{\mathit{Sh}}\nolimits}
\def\NL{\mathop{N\!L}\nolimits}
\def\CH{\mathop{\mathrm{CH}}\nolimits}
\def\proetale{{pro\text{-}\acute{e}tale}}
\def\etale{{\acute{e}tale}}
\def\QCoh{\mathit{QCoh}}
\def\Ker{\mathop{\mathrm{Ker}}}
\def\Im{\mathop{\mathrm{Im}}}
\def\Coker{\mathop{\mathrm{Coker}}}
\def\Coim{\mathop{\mathrm{Coim}}}

% Boxtimes
%
\DeclareMathSymbol{\boxtimes}{\mathbin}{AMSa}{"02}

%
% Macros for moduli stacks/spaces
%
\def\QCohstack{\mathcal{QC}\!\mathit{oh}}
\def\Cohstack{\mathcal{C}\!\mathit{oh}}
\def\Spacesstack{\mathcal{S}\!\mathit{paces}}
\def\Quotfunctor{\mathrm{Quot}}
\def\Hilbfunctor{\mathrm{Hilb}}
\def\Curvesstack{\mathcal{C}\!\mathit{urves}}
\def\Polarizedstack{\mathcal{P}\!\mathit{olarized}}
\def\Complexesstack{\mathcal{C}\!\mathit{omplexes}}
% \Pic is the operator that assigns to X its picard group, usage \Pic(X)
% \Picardstack_{X/B} denotes the Picard stack of X over B
% \Picardfunctor_{X/B} denotes the Picard functor of X over B
\def\Pic{\mathop{\mathrm{Pic}}\nolimits}
\def\Picardstack{\mathcal{P}\!\mathit{ic}}
\def\Picardfunctor{\mathrm{Pic}}
\def\Deformationcategory{\mathcal{D}\!\mathit{ef}}
\usepackage{cleveref}
\crefname{originallemma}{original lemma}{original lemmas}
\Crefname{originallemma}{Original lemma}{Original lemmas}
\makeatletter
\newlabel{derived-lemma-two-ss-complex-functor}{{015J}{source}{}{sourcepointer-1}{}}
\newlabel{derived-lemma-two-ss-complex-functor@cref}{{[originallemma][0][]015J}{[1][0][]source}{}{}{}}
\makeatother
\crefname{originalreference}{original reference}{original references}
\Crefname{originalreference}{Original reference}{Original references}
\makeatletter
\newlabel{lemma-colimit}{{00D5}{source}{}{sourcepointer-2}{}}
\newlabel{lemma-colimit@cref}{{[originalreference][0][]00D5}{[1][0][]source}{}{}{}}
\makeatother
\title{Cohomology of a torsion sheaf, and of a complex with torsion cohomology, on a quasi-compact quasi-separated scheme is torsion}
\author{The Stacks Project Authors}
\date{stacks-project:etale-cohomology.tex:local-snapshot-20261011; GNU FDL 1.2 or later}
\begin{document}
\maketitle
\noindent\textit{Editorial note.} All delivered mathematical text is the author's own, copied from the pinned Stacks Project chapter file etale-cohomology.tex; no statement or proof step was rewritten, reordered or removed.\par\medskip
\noindent\textit{Editorial note.} The counted proof is the author's complete argument: the proof is written out in full: the exhaustion of a torsion sheaf by its d-torsion subsheaves with the colimit argument, and the spectral sequence reduction for complexes.\par\medskip
\noindent\textit{Editorial note.} The label lemma-torsion-cohomology names the counted statement, and the proof environment immediately following it in the source is delivered verbatim as its proof.\par\medskip
\noindent\textit{Editorial note.} The counted statement is carried by the author's own lemma environment in the chapter source.\par\medskip
\noindent\textit{Editorial note.} The statement and proof are copied from the pinned chapter file; the two cited inputs are kept as external pointers with their canonical Stacks tags.\par\medskip
\noindent\textit{Allowed standard prerequisites (editorial).}\begin{itemize}
\item Commutation of cohomology with filtered colimits. Same chapter, the lemma computing the cohomology of a filtered colimit of sheaves as the colimit of the cohomologies. Supplies the identity H\textasciicircum{}n(X, F) = colim H\textasciicircum{}n(X, F[d]) used in part (1).
\item The spectral sequence of a complex computing its cohomology. Derived Categories, the lemma giving the spectral sequence E\_2\textasciicircum{}\{p,q\} = H\textasciicircum{}p(X, H\textasciicircum{}q(K)) converging to H\textasciicircum{}n(X, K). Supplies the reduction of part (2) to the case of sheaves.
\end{itemize}\medskip
\ifdefined\hypertarget\hypertarget{sourcepointer-1}{}\fi
\noindent\textit{Source reference.} Original lemma 015J: sources/stacks-chow/tags.txt, source lines 1510--1510. Prerequisite pointer invoked at proof line(s) 13919; it names a standard fact of a sibling chapter that the counted argument uses, and every step of the counted proof built on it is retained verbatim in the delivered body.\par\medskip
\ifdefined\hypertarget\hypertarget{sourcepointer-2}{}\fi
\noindent\textit{Source reference.} Original reference 00D5: sources/stacks-etale-cohomology/pinned-original-source.tex, source lines 7427--7427. Prerequisite pointer invoked at proof line(s) 13912; it names a standard fact of the same chapter that the counted argument uses, and every step of the counted proof built on it is retained verbatim in the delivered body.\par\medskip
\begin{lemma}
\label{lemma-torsion-cohomology}
Let $X$ be a quasi-compact and quasi-separated scheme.
\begin{enumerate}
\item If $\mathcal{F}$ is a torsion abelian sheaf on $X_\etale$, then
$H^n_\etale(X, \mathcal{F})$ is a torsion abelian group for all $n$.
\item If $K$ in $D^+(X_\etale)$ has torsion cohomology sheaves, then
$H^n_\etale(X, K)$ is a torsion abelian group for all $n$.
\end{enumerate}
\end{lemma}

\begin{proof}
To prove (1) we write $\mathcal{F} = \bigcup \mathcal{F}[n]$ where
$\mathcal{F}[d]$ is the $d$-torsion subsheaf. By
Lemma \ref{lemma-colimit} we have
$H^n_\etale(X, \mathcal{F}) = \colim H^n_\etale(X, \mathcal{F}[d])$.
This proves (1) as $H^n_\etale(X, \mathcal{F}[d])$ is annihilated by $d$.

\medskip\noindent
To prove (2) we can use the spectral sequence
$E_2^{p, q} = H^p_\etale(X, H^q(K))$ converging to $H^n_\etale(X, K)$
(Derived Categories, Lemma \ref{derived-lemma-two-ss-complex-functor})
and the result for sheaves.
\end{proof}

\end{document}
